% TCC · Bacharelado em Ciência da Computação · UNITRI · 2026/2
% Modelo abnTeX2 (classe abntex2). No Overleaf, compilar com pdfLaTeX.
%
% Marcadores usados no texto:
%   \pendente{...}  trecho a escrever ou a confirmar (sai em vermelho no PDF)
%   \resultado      número que só pode vir de execução real do modelo ou do app
% Nenhum dos dois pode sobrar na versão postada.

\documentclass[
  12pt,
  oneside,
  a4paper,
  chapter=TITLE,
  section=TITLE,
  english,
  brazil
]{abntex2}

\usepackage{lmodern}
\usepackage[T1]{fontenc}
\usepackage[utf8]{inputenc}
\usepackage{indentfirst}
\usepackage{xcolor}
\usepackage{graphicx}
\usepackage{microtype}
\usepackage{lastpage}
\usepackage{tabularx}
% Citações autor-data (NBR 10520). O grupo decidiu pela versão de 2023 (autor em
% maiúsculas e minúsculas); este pacote segue a de 2002 (autor em caixa alta). Falta ajustar.
% abnt-emphasize=bf: título da obra em negrito nas referências, como o manual recomenda.
\usepackage[alf, abnt-emphasize=bf]{abntex2cite}

\graphicspath{{figuras/}}

\newcommand{\pendente}[1]{\textcolor{red}{[#1]}}
\newcommand{\resultado}{\textcolor{red}{[RESULTADO A PREENCHER]}}

% Título provisório
\titulo{Classificação da maturação do tomate em dispositivos móveis: um aplicativo offline de apoio à decisão de colheita}
\autor{Leonardo Sergi Molina \\ Vitor Kazuma \pendente{SOBRENOME} \\ Victor José \pendente{SOBRENOME}}
\local{Uberlândia}
\data{2026}
\orientador{\pendente{NOME DO ORIENTADOR}}
\instituicao{%
  Centro Universitário do Triângulo -- UNITRI
  \par
  Bacharelado em Ciência da Computação}
\tipotrabalho{Monografia (Bacharelado)}
\preambulo{Monografia apresentada ao curso de Bacharelado em Ciência da Computação do Centro Universitário do Triângulo (UNITRI), como requisito parcial para a obtenção do título de Bacharel em Ciência da Computação.}

\hypersetup{
  pdftitle={Classificação da maturação do tomate em dispositivos móveis},
  pdfauthor={Grupo do TCC},
  pdfsubject={Monografia de Bacharelado em Ciência da Computação},
  pdfkeywords={tomate}{maturação}{YOLO}{aplicativo móvel},
  colorlinks=true,
  linkcolor=blue,
  citecolor=blue,
  urlcolor=blue,
  bookmarksdepth=4
}

% Margens e espaçamento (NBR 14724)
\setlrmarginsandblock{3cm}{2cm}{*}
\setulmarginsandblock{3cm}{2cm}{*}
\checkandfixthelayout
\setlength{\parindent}{1.3cm}
\setlength{\parskip}{0.2cm}

\begin{document}

\selectlanguage{brazil}
\frenchspacing

% ---------------------------------------------------------------- pré-textuais
\imprimircapa
\imprimirfolhaderosto*

% Resumo (NBR 6028): parágrafo único, de 150 a 500 palavras, escrito por último.
\setlength{\absparsep}{18pt}
\begin{resumo}
\pendente{Escrever por último: objetivo, método, principais resultados e conclusão, em um único parágrafo.}

\vspace{\onelineskip}

\noindent
\textbf{Palavras-chave}: maturação do tomate. visão computacional. YOLO. aplicativo móvel. agricultura digital.
\end{resumo}

% Abstract: versão em inglês do resumo, revisada por uma pessoa.
\begin{resumo}[Abstract]
\begin{otherlanguage*}{english}
\pendente{Write last: English version of the Portuguese abstract.}

\vspace{\onelineskip}

\noindent
\textbf{Keywords}: tomato ripeness. computer vision. YOLO. mobile application. digital agriculture.
\end{otherlanguage*}
\end{resumo}


\pdfbookmark[0]{\listfigurename}{lof}
\listoffigures*
\cleardoublepage

\pdfbookmark[0]{\listtablename}{lot}
\listoftables*
\cleardoublepage

% Lista de siglas: manter em ordem alfabética e só com siglas usadas no texto.
\begin{siglas}
  \item[ABNT] Associação Brasileira de Normas Técnicas
  \item[APK] Android Package Kit
  \item[CNN] Convolutional Neural Network (rede neural convolucional)
  \item[EXIF] Exchangeable Image File Format
  \item[IA] Inteligência Artificial
  \item[IBGE] Instituto Brasileiro de Geografia e Estatística
  \item[LLM] Large Language Model (modelo de linguagem de grande escala)
  \item[MER] Modelo Entidade-Relacionamento
  \item[MVP] Minimum Viable Product (produto mínimo viável)
  \item[UML] Unified Modeling Language
  \item[YOLO] You Only Look Once
\end{siglas}


\pdfbookmark[0]{\contentsname}{toc}
\tableofcontents*
\cleardoublepage

% ---------------------------------------------------------------- textuais
\textual

\chapter{Introdução}\label{cap:introducao}

% Rascunho inicial (04/11/2026). O manual pede que a Introdução seja reescrita
% por último, para refletir o trabalho concluído. Todo parágrafo de contexto,
% problema e justificativa precisa de citação.

\section{Contextualização}

A agricultura familiar responde pela maior parte dos estabelecimentos rurais do Brasil. Segundo o Censo Agropecuário 2017, são 3.897.408 estabelecimentos familiares, o equivalente a 76,8\% do total. Na produção de tomate estaqueado, a participação é ainda maior: 82,2\% dos estabelecimentos são familiares, e Minas Gerais responde por 24,0\% da produção nacional \cite{ibge2019censo}.

Parte relevante dos alimentos se perde antes de chegar ao consumidor. A FAO estima que 13,8\% dos alimentos produzidos em 2016 se perderam entre a pós-colheita e o varejo, sem contar o próprio varejo, e aponta o estádio de maturação e o momento da colheita entre as causas de perda relatadas na colheita de frutas \cite{fao2019}.

O uso de tecnologias digitais já é comum no campo: em pesquisa com 504 produtores brasileiros, 84\% declararam usar ao menos uma tecnologia digital \cite{bolfe2020}. A conectividade, porém, segue limitada. A rede móvel funciona em 68,0\% dos domicílios rurais, contra 96,1\% dos urbanos \cite{ibge2026pnadtic}, o que torna arriscado depender de internet numa ferramenta pensada para uso na lavoura.

Técnicas de visão computacional baseadas em aprendizado profundo já foram aplicadas à classificação da maturação de frutos, inclusive em aplicativos para smartphone \cite{rizzo2023,mohamedon2021}. No Brasil, \citeonline{mundimfilho2021} usaram uma rede neural convolucional para indicar se cafeeiros estavam prontos para a colheita e obtiveram acurácia de 93\%.

\section{Problema de pesquisa}

Diante desse cenário, este trabalho busca responder à seguinte pergunta: em que medida um classificador leve, executado no próprio celular e sem conexão com a internet, consegue indicar se o tomate está no ponto de colheita a partir de uma foto tirada no campo?

\section{Justificativa}

\pendente{Relevância para o produtor (decisão de colheita e redução de perdas), para a área acadêmica (avaliação de um modelo embarcado com fotos de campo) e diferença em relação a \citeonline{mundimfilho2021}: outra cultura, modelo recente, avaliação com fotos próprias e aplicativo multiplataforma. Todo parágrafo com citação.}

\section{Objetivos}

\subsection{Objetivo geral}

Desenvolver um aplicativo móvel que, a partir de uma foto, classifique a maturação do tomate e recomende colher ou esperar, funcionando sem conexão com a internet.

\subsection{Objetivos específicos}

\begin{alineas}
  \item definir as classes de maturação a partir de escalas de cor de referência;
  \item treinar e avaliar um classificador leve com um conjunto de dados público e fotos próprias de teste;
  \item desenvolver o aplicativo com inferência no dispositivo, regra de decisão e histórico de análises;
  \item avaliar o desempenho no celular quanto à acurácia, ao F1 por classe, ao tempo de resposta e ao funcionamento sem internet.
\end{alineas}

\section{Metodologia}

\pendente{Resumo de um parágrafo: pesquisa aplicada, quantitativa e experimental; etapas de dados, treinamento, desenvolvimento e avaliação. Detalhes no Capítulo~\ref{cap:metodologia}.}

\section{Hipótese}

Parte-se da hipótese de que um classificador embarcado atinge acurácia de pelo menos \pendente{META A DEFINIR COM O PROFESSOR} no conjunto de teste com fotos próprias, responde em poucos segundos e funciona sem internet.

\section{Estrutura do trabalho}

Além desta introdução, o trabalho está organizado em cinco capítulos. O Capítulo~\ref{cap:referencial} apresenta os conceitos de maturação do tomate, visão computacional, aprendizado profundo e computação móvel que fundamentam a proposta. O Capítulo~\ref{cap:estado-da-arte} analisa trabalhos correlatos e posiciona esta pesquisa em relação a eles. O Capítulo~\ref{cap:metodologia} descreve os dados, o treinamento do modelo, os requisitos, a modelagem do sistema e o plano de avaliação. O Capítulo~\ref{cap:resultados} apresenta o aplicativo e os resultados obtidos, e o Capítulo~\ref{cap:conclusao} traz as conclusões, as limitações e as sugestões de trabalhos futuros.

\chapter{Referencial teórico}\label{cap:referencial}

% Conceitos que sustentam o trabalho. Fontes: livros, artigos e teses (nada de
% Wikipédia, blogs ou notas de aula). Comparar autores em vez de listar resumos.
% As chaves entre colchetes são entradas do referencias.bib.

\pendente{Parágrafo de abertura: o que o capítulo apresenta e em que ordem.}

\section{Maturação do tomate}\label{sec:maturacao}

% Estádios de maturação e fisiologia: seymour2013. Fruto climatérico (continua
% amadurecendo depois de colhido): flores2011. Escalas de cor: usda1991
% (6 classes por % da superfície colorida), brasil1995 (5 subgrupos; revogada pela
% brasil2018, citar como escala histórica e de mercado) e pbmh2003. Ponto de
% colheita conforme o destino da venda: pbmh2003.
\pendente{A escrever.}

\section{Processamento digital de imagens e visão computacional}

% Imagem digital, espaços de cor (RGB, HSV, Lab) e por que a cor indica a maturação;
% limites da abordagem clássica por cor (iluminação, fundo): rizzo2023.
\pendente{A escrever.}

\section{Aprendizado profundo e redes neurais convolucionais}

% lecun2015, goodfellow2016.
\pendente{A escrever.}

\section{Classificação de imagens e a família YOLO}

% Diferença entre classificação e detecção. YOLO: redmon2016 (original),
% terven2023 (revisão até o YOLOv8), badgujar2024 (YOLO na agricultura),
% jocher2026yolo26 (YOLO26, preprint) e jocher2023ultralytics (software).
\pendente{A escrever.}

\section{Transfer learning}

% pan2010; redes leves para celular: howard2017, sandler2018, tan2019.
\pendente{A escrever.}

\section{Métricas de avaliação}

% Acurácia, precisão, recall, F1 e matriz de confusão: sokolova2009.
% Contexto de benchmarks de visão: everingham2010, lin2014.
\pendente{A escrever.}

\section{Computação móvel e inferência no dispositivo}

% Inferência no aparelho e na nuvem; formatos LiteRT e Core ML; plugin
% ultralytics2026flutter; escolha do framework: oliveira2023.
% Contexto de uso: statcounter2026 (Android no Brasil), ibge2026pnadtic e
% cetic2025 (conectividade rural).
\pendente{A escrever.}

\chapter{Estado da arte}\label{cap:estado-da-arte}

% Manual: 4 a 5 trabalhos dos últimos 5 a 10 anos. Para cada um, responder:
% qual problema tentou resolver, como abordou (algoritmos, técnicas, arquitetura),
% quais foram os principais achados e o que faltou. Depois, comparar os trabalhos
% entre si e com este TCC, fechar com o quadro comparativo e deixar clara a lacuna.
% Ler cada artigo antes de escrever e conferir os números no PDF.

\pendente{Parágrafo de abertura: critérios de busca (bases, período, palavras-chave) e de seleção dos trabalhos.}

\section{Decisão de colheita do café com rede neural convolucional}

% mundimfilho2021 (artigo do SBSI) e santos2022 (dissertação da mesma linha, UFU).
% Antecedente mais próximo: aplicativo Android offline, em português, que responde
% "colher" ou "não colher". Deixar explícito em que este TCC se diferencia.
\pendente{A escrever.}

\section{Detecção de frutos de café em ramos com YOLO}

% bazame2023: fotos de smartphone em Patos de Minas-MG; comparação entre
% versões do YOLO e efeito do tamanho da imagem de entrada.
\pendente{A escrever.}

\section{Aplicativo de maturação de azeitonas com YOLOv8}

% alvino2024: aplicativo brasileiro (Flutter, API em Flask e YOLOv8) que depende de servidor.
\pendente{A escrever.}

\section{Detecção leve da maturação do tomate}

% dong2025: GPC-YOLO, tomate em ambiente não estruturado, 3 classes; não foi testado em celular.
\pendente{A escrever.}

\section{Aplicativo Android com YOLO leve executado no dispositivo}

% zheng2025: arroz; modelo embarcado com testes de desempenho no próprio smartphone.
% Bom modelo de metodologia para a avaliação no celular.
\pendente{A escrever.}

\section{Quadro comparativo e posicionamento deste trabalho}

\pendente{Parágrafo comparando as abordagens (diálogo entre os autores) e deixando clara a lacuna que este trabalho preenche.}

\begin{table}[htb]
  \ABNTEXfontereduzida
  \caption{Comparação entre os trabalhos correlatos e este trabalho}
  \label{tab:comparativo}
  \begin{tabularx}{\textwidth}{lXXXX}
    \hline
    \textbf{Trabalho} & \textbf{Cultura} & \textbf{Modelo} & \textbf{Onde roda} & \textbf{Saída para o produtor} \\
    \hline
    \citeonline{mundimfilho2021} & Café & CNN (MobileNet) & \pendente{conferir} & Colher ou não colher \\
    \citeonline{bazame2023} & Café & YOLOv3 e YOLOv4 & \pendente{conferir} & \pendente{conferir} \\
    \citeonline{alvino2024} & Azeitona & YOLOv8 & Servidor & Índice de maturação \\
    \citeonline{dong2025} & Tomate & GPC-YOLO (YOLOv8n modificado) & \pendente{conferir} & \pendente{conferir} \\
    \citeonline{zheng2025} & Arroz & YOLO leve & No smartphone & \pendente{conferir} \\
    Este trabalho & Tomate & YOLO26n-cls & No celular, offline & Colher ou esperar \\
    \hline
  \end{tabularx}
  \legend{Fonte: elaborada pelos autores.}
\end{table}

\chapter{Metodologia}\label{cap:metodologia}

% Manual (trabalho prático): ferramentas usadas, critérios de escolha e por que
% não outras; os 4 diagramas UML (casos de uso, classes, sequência e atividades)
% e o MER, cada um descrito; público-alvo definido. Texto no passado e na voz
% passiva ("foi realizado", "optou-se por").

\section{Classificação da pesquisa}

\pendente{Pesquisa aplicada, de abordagem quantitativa e procedimento experimental. Justificar cada classificação.}

\section{Público-alvo}

\pendente{Produtores de tomate, sobretudo pequenos e familiares, que precisam decidir o momento da colheita no campo, muitas vezes sem internet.}

\section{Etapas do trabalho}

\pendente{Fluxograma das etapas: escolha do dataset, definição das classes, treinamento, desenvolvimento do aplicativo, testes e avaliação.}

\section{Ferramentas e justificativas}

\pendente{Para cada ferramenta (Flutter, pacote Ultralytics e YOLO26n-cls, LiteRT e Core ML, SQLite, Google Colab, Overleaf, GitHub), explicar o critério da escolha e por que outras opções não foram usadas.}

\section{Conjunto de dados e classes}\label{sec:dados}

\pendente{Dataset público usado (origem, licença, número de imagens), mapeamento das classes originais para verde, de vez, maduro, passado e outro, e como foram gerados os recortes. Conjunto de teste próprio: tomates comprados verdes e fotografados diariamente até amadurecerem, com 2 ou 3 celulares e variação de luz. Separação entre treino e teste por fruto ou lote, nunca por foto.}

\section{Treinamento do modelo}

\pendente{Versão fixada do pacote, pesos pré-treinados, tamanho de imagem, número de épocas, hardware e semente. Tudo deve ser reprodutível a partir do notebook do repositório.}

\section{Requisitos}

\begin{table}[htb]
  \ABNTEXfontereduzida
  \caption{Requisitos funcionais e não funcionais do aplicativo}
  \label{tab:requisitos}
  \begin{tabularx}{\textwidth}{lX}
    \hline
    \textbf{Código} & \textbf{Requisito} \\
    \hline
    RF01 & Capturar uma foto pela câmera ou escolher uma imagem da galeria. \\
    RF02 & Classificar a maturação do tomate no próprio dispositivo, sem conexão com a internet. \\
    RF03 & Exibir a classe, a confiança e a recomendação de colher ou esperar. \\
    RF04 & Alertar o usuário quando a foto for inadequada ou a confiança for baixa. \\
    RF05 & Armazenar cada análise no histórico do dispositivo. \\
    RF06 & Consultar o histórico de análises. \\
    RNF01 & Funcionar integralmente sem conexão com a internet. \\
    RNF02 & Executar em Android 7.0 ou superior. \\
    RNF03 & Responder em até \pendente{META DE TEMPO} por análise. \\
    RNF04 & Não enviar fotos nem dados de localização para servidores. \\
    RNF05 & Manter a interface legível sob sol forte, com alto contraste e botões grandes. \\
    \hline
  \end{tabularx}
  \legend{Fonte: elaborada pelos autores.}
\end{table}

\section{Modelagem do sistema}

\subsection{Diagrama de casos de uso}
\pendente{Figura e descrição: fotografar fruto, analisar maturação, ver recomendação e consultar histórico. Ator: produtor.}

\subsection{Diagrama de classes}
\pendente{Figura e descrição: Analise, Cultura, ClasseMaturacao, RegraDecisao, ClassificadorLocal e RepositorioAnalises.}

\subsection{Diagrama de sequência}
\pendente{Figura e descrição: foto, pré-processamento, classificador, regra de decisão, tela de resultado e histórico.}

\subsection{Diagrama de atividades}
\pendente{Figura e descrição, com as decisões: a foto está boa? a confiança é suficiente? a recomendação é colher ou esperar?}

\section{Modelo entidade-relacionamento}

\pendente{Figura e descrição: Cultura (1:N) ClasseMaturacao (1:1) RegraDecisao; Analise (N:1) ClasseMaturacao, com data, caminho da foto no dispositivo, confiança e versão do modelo.}

\section{Plano de testes e de avaliação}

\pendente{Casos de teste derivados dos requisitos; métricas do modelo (acurácia, precisão, recall, F1 por classe e matriz de confusão) no conjunto de teste próprio; tempo de resposta medido em 2 ou 3 celulares; funcionamento em modo avião.}

\section{Declaração de uso de inteligência artificial}

Durante a elaboração deste trabalho, foram utilizadas ferramentas de inteligência artificial generativa como apoio à pesquisa bibliográfica, à escrita e revisão do texto e ao desenvolvimento do código. Todo o conteúdo produzido com esse apoio foi revisado pelos autores, as referências citadas foram conferidas nas fontes originais e todos os resultados apresentados vêm de execuções reais do modelo e do aplicativo. \pendente{Listar as ferramentas usadas, com nome e fabricante.}

\chapter{Resultados e discussão}\label{cap:resultados}

% Manual: não usar "Desenvolvimento" como título. Apresentar o software com prints,
% tabelas com título acima e fonte abaixo, toda figura citada no texto. Comparar
% com o Estado da Arte e ser honesto sobre as dificuldades.
% Todo número deste capítulo sai de execução real: até lá, usar \resultado.

\section{Desempenho do modelo}

\pendente{Tabela com acurácia, precisão, recall e F1 por classe no conjunto de teste próprio (\resultado), matriz de confusão e discussão das classes que mais se confundem.}

\section{O aplicativo}

\pendente{Uma subseção por tela (câmera, resultado, histórico), com print e descrição de cada função.}

\section{Testes funcionais e de desempenho}

\pendente{Resultado dos casos de teste, tempo de resposta por aparelho (\resultado) e funcionamento em modo avião.}

\section{Comparação com os trabalhos correlatos}

\pendente{Retomar o quadro do Capítulo~\ref{cap:estado-da-arte}: onde este trabalho avança, confirma ou diverge.}

\section{Dificuldades e limitações}

\pendente{Diferença entre as fotos do dataset e as de campo, subjetividade da classe ``de vez'', tamanho do conjunto de teste, entre outras.}

\chapter{Conclusão}\label{cap:conclusao}

% Manual: responder à pergunta de pesquisa, dizer se os objetivos e a hipótese
% foram atingidos, destacar contribuições e limitações e sugerir trabalhos futuros.
% Sem dados novos e, em regra, sem citações. Texto corrido, em sintonia com a
% Introdução.

\pendente{Resposta ao problema de pesquisa e verificação dos objetivos e da hipótese.}

\pendente{Contribuições e limitações.}

\pendente{Trabalhos futuros: faixa estimada de dias até a colheita, comparação com um LLM multimodal na nuvem e outras culturas (se não tiverem entrado no trabalho).}


% ---------------------------------------------------------------- pós-textuais
\postextual

\bibliography{referencias}

\begin{apendicesenv}
\partapendices
\chapter{Regra de decisão usada pelo aplicativo}\label{ap:regra}

A Tabela~\ref{tab:regra} relaciona cada classe do modelo às escalas de cor de referência e à recomendação exibida ao produtor.

\begin{table}[htb]
  \ABNTEXfontereduzida
  \caption{Classes do modelo, escalas de referência e recomendação}
  \label{tab:regra}
  \begin{tabularx}{\textwidth}{lXXl X}
    \hline
    \textbf{Classe} & \textbf{Portaria 553/1995} & \textbf{USDA} & \textbf{Vermelho} & \textbf{Recomendação} \\
    \hline
    Verde & Todo verde e verde maduro & Green e Breakers & até 10\% & Esperar \\
    De vez & Pintado e rosado & Turning e Pink & 10\% a 60\% & Colher se for para mercado distante; senão, esperar \\
    Maduro & Vermelho e vermelho maduro & Light red e Red & acima de 60\% & Colher \\
    Passado & Defeito grave & -- & -- & Colher já ou descartar \\
    Outro & -- & -- & -- & Pedir nova foto \\
    \hline
  \end{tabularx}
  \legend{Fonte: elaborada pelos autores com base em \citeonline{brasil1995}, \citeonline{usda1991} e \citeonline{pbmh2003}.}
\end{table}

\end{apendicesenv}

\end{document}
